% GENERATED FILE. Edit canonical lessons, then run npm run generate:books.
% canonical-source-hash: fnv1a64:beeaf519ffbab9e1
% canonical-lessons: IT-C415-noleggiare, IT-C415-noleggio, IT-C415-targa, IT-C415-libretto, IT-C415-assicurazione

\chapter{A noleggio}
\label{ch:it-a-noleggio}
\hlchaptermodality{415}

\begin{chapteropening}
\textbf{By the end of this chapter:} \emph{I can hire a car: hiring, the rental, the number plate, the logbook and the insurance.}

Complete the last lesson of chapter 415: I can hire a car: hiring, the rental, the number plate, the logbook and the insurance.
\end{chapteropening}

\section[noleggiare]{noleggiare — to hire, to rent (a car or a bike)}

\label{lesson:IT-C415-noleggiare}



\begin{quote}
Before the new one: say the Italian for the pothole, then the Italian for the dirt road.
\end{quote}

\subsection*{You'll want to know: noleggiare}
\textbf{noleggiare} — \textquotedblleft{}to hire, to rent (a car or a bike)\textquotedblright{}.

\begin{quote}
In vacanza noleggiamo una macchina.
\end{quote}

\begin{quote}
\emph{On holiday we hire a car.}
\end{quote}

For a car, a bike or a boat; for a flat, \emph{affittare}.

\begin{cousinweb}
\emph{noleggiare} is built on \emph{nolo} 'hire, freight charge', which is usually traced through Medieval Latin \emph{naulum} to Greek \emph{naûlon} 'the fare for a passage by ship'; it came into Italian as a borrowing.
\end{cousinweb}

\begin{grammarlens}[title={What you've built}]
One more word for telling, explaining and describing at B1.
\end{grammarlens}

\subsection*{Your turn}
\begin{itemize}
\raggedright
  \item \emph{Say it:} \emph{noleggiare}
  \item \emph{Say it:} \emph{noleggiare}, once more
  \item \emph{Say it:} a sentence of your own with \emph{noleggiare}
  \item \emph{Recall:} say the Italian for the pothole, then the Italian for the dirt road, then say \emph{noleggiare} again
\end{itemize}

\begin{tcolorbox}[breakable,colback=teal!4,colframe=teal!35!black,title={Before you move on}]
What does noleggiare mean? (\textquotedblleft{}To hire, to rent (a car or a bike)\textquotedblright{}.) Say it once more.
\end{tcolorbox}
\section[il noleggio]{il noleggio — the hire, the rental}

\label{lesson:IT-C415-noleggio}



\begin{quote}
Before the new one: say the Italian for the dirt road, then the Italian for to hire.
\end{quote}

\subsection*{You'll want to know: il noleggio}
\textbf{il noleggio} — \textquotedblleft{}the hire, the rental\textquotedblright{}.

\begin{quote}
Il noleggio costa quaranta euro al giorno.
\end{quote}

\begin{quote}
\emph{The hire costs forty euros a day.}
\end{quote}

\emph{Un'auto a noleggio} — a hire car; \emph{il noleggio bici} — bike hire.

\begin{cousinweb}
\emph{noleggio} is built on \emph{nolo} 'hire', which is usually traced through Medieval Latin \emph{naulum} to Greek \emph{naûlon} 'ship fare'; it came into Italian as a borrowing.
\end{cousinweb}

\begin{grammarlens}[title={What you've built}]
One more word for telling, explaining and describing at B1.
\end{grammarlens}

\subsection*{Your turn}
\begin{itemize}
\raggedright
  \item \emph{Say it:} \emph{il noleggio}
  \item \emph{Say it:} \emph{il noleggio}, once more
  \item \emph{Say it:} a sentence of your own with \emph{il noleggio}
  \item \emph{Recall:} say the Italian for the dirt road, then the Italian for to hire, then say \emph{il noleggio} again
\end{itemize}

\begin{tcolorbox}[breakable,colback=teal!4,colframe=teal!35!black,title={Before you move on}]
What does il noleggio mean? (\textquotedblleft{}The hire, the rental\textquotedblright{}.) Say it once more.
\end{tcolorbox}
\section[la targa]{la targa — the number plate; the plaque}

\label{lesson:IT-C415-targa}



\begin{quote}
Before the new one: say the Italian for to hire, then the Italian for the hire.
\end{quote}

\subsection*{You'll want to know: la targa}
\textbf{la targa} — \textquotedblleft{}the number plate; the plaque\textquotedblright{}.

\begin{quote}
Hai visto la targa di quella macchina?
\end{quote}

\begin{quote}
\emph{Did you see that car's number plate?}
\end{quote}

Also a plaque on a wall, in memory of someone.

\begin{cousinweb}
\emph{targa} first meant a shield. It is usually traced to a Germanic word for a shield, but the exact path is not settled.
\end{cousinweb}

\begin{grammarlens}[title={What you've built}]
One more word for telling, explaining and describing at B1.
\end{grammarlens}

\subsection*{Your turn}
\begin{itemize}
\raggedright
  \item \emph{Say it:} \emph{la targa}
  \item \emph{Say it:} \emph{la targa}, once more
  \item \emph{Say it:} a sentence of your own with \emph{la targa}
  \item \emph{Recall:} say the Italian for to hire, then the Italian for the hire, then say \emph{la targa} again
\end{itemize}

\begin{tcolorbox}[breakable,colback=teal!4,colframe=teal!35!black,title={Before you move on}]
What does la targa mean? (\textquotedblleft{}The number plate; the plaque\textquotedblright{}.) Say it once more.
\end{tcolorbox}
\section[il libretto]{il libretto — the logbook (of a car); the booklet}

\label{lesson:IT-C415-libretto}



\begin{quote}
Before the new one: say the Italian for the hire, then the Italian for the number plate.
\end{quote}

\subsection*{You'll want to know: il libretto}
\textbf{il libretto} — \textquotedblleft{}the logbook (of a car); the booklet\textquotedblright{}.

\begin{quote}
Il vigile chiede patente e libretto.
\end{quote}

\begin{quote}
\emph{The warden asks for the licence and the logbook.}
\end{quote}

\emph{Il libretto di circolazione} is the car's registration document; at university, \emph{il libretto} records your exams.

\begin{cousinweb}
\emph{libretto} is a little \emph{libro} 'book', which is inherited from Latin \emph{liber} 'book' (EtymDB).
\end{cousinweb}

\begin{grammarlens}[title={What you've built}]
One more word for telling, explaining and describing at B1.
\end{grammarlens}

\subsection*{Your turn}
\begin{itemize}
\raggedright
  \item \emph{Say it:} \emph{il libretto}
  \item \emph{Say it:} \emph{il libretto}, once more
  \item \emph{Say it:} a sentence of your own with \emph{il libretto}
  \item \emph{Recall:} say the Italian for the hire, then the Italian for the number plate, then say \emph{il libretto} again
\end{itemize}

\begin{tcolorbox}[breakable,colback=teal!4,colframe=teal!35!black,title={Before you move on}]
What does il libretto mean? (\textquotedblleft{}The logbook (of a car); the booklet\textquotedblright{}.) Say it once more.
\end{tcolorbox}
\section[l'assicurazione]{l'assicurazione — the insurance}

\label{lesson:IT-C415-assicurazione}



\begin{quote}
Before the new one: say the Italian for the number plate, then the Italian for the logbook.
\end{quote}

\subsection*{You'll want to know: l'assicurazione}
\textbf{l'assicurazione} — \textquotedblleft{}the insurance\textquotedblright{}.

\begin{quote}
L'assicurazione della macchina scade a maggio.
\end{quote}

\begin{quote}
\emph{The car insurance runs out in May.}
\end{quote}

\emph{Assicurare} — to insure, and to assure.

\begin{cousinweb}
\emph{assicurazione} is built on \emph{assicurare} 'to make safe', from \emph{a-} and \emph{sicuro} 'safe', which is inherited from Latin \emph{sēcūrus} 'free from care' (Lewis \& Short).
\end{cousinweb}

\begin{grammarlens}[title={What you've built}]
One more word for telling, explaining and describing at B1.
\end{grammarlens}

\subsection*{Your turn}
\begin{itemize}
\raggedright
  \item \emph{Say it:} \emph{l'assicurazione}
  \item \emph{Say it:} \emph{l'assicurazione}, once more
  \item \emph{Say it:} a sentence of your own with \emph{l'assicurazione}
  \item \emph{Recall:} say the Italian for the number plate, then the Italian for the logbook, then say \emph{l'assicurazione} again
\end{itemize}

\begin{tcolorbox}[breakable,colback=teal!4,colframe=teal!35!black,title={Before you move on}]
What does l'assicurazione mean? (\textquotedblleft{}The insurance\textquotedblright{}.) Say it once more.
\end{tcolorbox}
